\documentclass[a4paper,12pt]{article}

\usepackage[T2A]{fontenc}
\usepackage[utf8]{inputenc}
\usepackage[russian]{babel}
\usepackage{amsmath,amssymb,mathtools}
\usepackage{PTSans}
\renewcommand{\familydefault}{\sfdefault}
\usepackage{icomma}
\usepackage[a4paper,margin=18mm,top=19mm,bottom=20mm]{geometry}
\usepackage{microtype}
\usepackage{tikz}
\usetikzlibrary{calc,arrows.meta,positioning}
\usepackage{tkz-euclide}
\usepackage{pgfplots}
\pgfplotsset{compat=1.18}
\usepackage{tabularx,booktabs,array}
\usepackage{enumitem}
\usepackage{hyperref}
\usepackage{parskip}
\usepackage{fancyhdr}
\usepackage{setspace}
\usepackage{xcolor}

\definecolor{accent}{HTML}{176B73}
\definecolor{darkaccent}{HTML}{124A50}
\definecolor{textgray}{HTML}{30383B}
\definecolor{softgray}{HTML}{69767A}
\definecolor{warm}{HTML}{B66A3C}

\hypersetup{colorlinks=true,linkcolor=darkaccent,urlcolor=darkaccent}
\setlength{\parindent}{0pt}
\setlength{\parskip}{6pt}
\onehalfspacing
\setlist[itemize]{leftmargin=6mm,itemsep=2pt,topsep=3pt}
\setlist[enumerate]{leftmargin=8mm,itemsep=5pt,topsep=3pt}
\renewcommand{\arraystretch}{1.25}
\renewcommand{\today}{4 октября 2026 г.}

\pagestyle{fancy}
\fancyhf{}
\fancyfoot[C]{\small\color{softgray}\thepage}
\renewcommand{\headrulewidth}{0pt}

\newcommand{\sectiontitle}[1]{%
  \vspace{2mm}%
  {\Large\bfseries\color{darkaccent}#1}\par\vspace{2mm}%
  {\color{accent!55}\rule{\linewidth}{0.6pt}}\par\vspace{2mm}%
}
\newcommand{\keyidea}[1]{\textbf{\color{darkaccent}#1}}
\newcommand{\vect}[2]{\overrightarrow{#1#2}}

\begin{document}

% ---------------- TITLE ----------------
\thispagestyle{empty}
\begin{center}
\vspace*{28mm}
{\small\color{softgray}МАТЕМАТИКА • ОГЭ}\par
\vspace{9mm}
{\fontsize{27}{31}\selectfont\bfseries\color{darkaccent}Чек-лист}\par
\vspace{4mm}
{\LARGE\bfseries Математическая модель, график и векторный маршрут}\par
\vspace{8mm}
{\large Движение по кругу • модуль и параметр • кусочные функции • векторы}\par
\vspace{21mm}
{\large София Кальней}\par
\vspace{4mm}
{\normalsize 4 октября 2026 г.}\par
\vfill
{\small\color{textgray}Лёвин Артём Александрович эксклюзивно для Софии Кальней}\par
\vspace{22mm}
\end{center}

\begin{tikzpicture}[remember picture,overlay]
  \draw[accent!55,line width=1.25pt]
    ([xshift=8mm,yshift=20mm]current page.south west)
    .. controls +(45mm,20mm) and +(-55mm,-5mm) ..
    ([xshift=-8mm,yshift=28mm]current page.south east);
  \draw[warm!45,line width=0.7pt]
    ([xshift=16mm,yshift=14mm]current page.south west)
    .. controls +(55mm,17mm) and +(-65mm,-8mm) ..
    ([xshift=-16mm,yshift=20mm]current page.south east);
\end{tikzpicture}

\clearpage

% ---------------- TEXT PROBLEMS ----------------
\sectiontitle{1. Движение по кругу: сначала выберите правильный объект поиска}

\keyidea{Главная идея.} Если задача спрашивает время одного круга, удобно именно это время обозначить через $x$. Если спрашивается скорость --- обозначайте скорость. Такой выбор часто убирает лишние неизвестные и делает уравнение коротким.

\textbf{Базовая связь:}
\[
S=vt,\qquad v=\frac{S}{t},\qquad t=\frac{S}{v}.
\]

\textbf{Для круговой трассы особенно полезны две модели:}
\begin{itemize}
  \item обычные единицы: км/ч, м/с, минуты, километры;
  \item \emph{круги в минуту}: если один круг проходится за $x$ минут, то скорость равна $\dfrac{1}{x}$ круга в минуту.
\end{itemize}

\textbf{Чек-лист перед уравнением:}
\begin{enumerate}
  \item Что требуется найти? Можно ли именно это обозначить через $x$?
  \item Все ли времена приведены к одним единицам?
  \item Что означает «опережает на один круг»? Это разность путей, равная длине круга.
  \item Что означает «осталось $400$ м до конца круга»? Пройденный путь меньше длины круга на $0{,}4$ км.
  \item Можно ли сократить неизвестную длину круга, разделив уравнение на неё?
\end{enumerate}

\subsection*{Показательный пример: лыжники}
Второй лыжник проходит круг за $x$ минут, первый --- на $2$ минуты быстрее, то есть за $x-2$ минут. За час первый опережает второго ровно на один круг.

Если длина круга равна $L$, то за $60$ минут:
\[
\frac{60L}{x-2}-\frac{60L}{x}=L.
\]
Делим на $L\ne0$:
\[
\frac{60}{x-2}-\frac{60}{x}=1,\qquad x>2.
\]
Отсюда
\[
\frac{120}{x(x-2)}=1
\quad\Longrightarrow\quad
x^2-2x-120=0
\quad\Longrightarrow\quad
x=12.
\]

\keyidea{Вывод:} второй лыжник проходит круг за $12$ минут. Здесь длина круга вообще не потребовалась.

\subsection*{Типичные ошибки}
\begin{itemize}
  \item обозначить через $x$ величину, которую затем приходится долго исключать;
  \item смешать минуты и часы в одном уравнении;
  \item забыть, что разность путей при опережении на круг равна ровно длине круга;
  \item получить отрицательное или слишком малое время и не проверить смысл ограничения $x>d$, где $d$ --- разница во времени прохождения круга.
\end{itemize}

% ---------------- MODULE GRAPH ----------------
\sectiontitle{2. График с модулем: сначала раскрыть по случаям}

Рассмотрим функцию занятия
\[
f(x)=|x|(x-1)-2x.
\]

Модуль раскрывается по знаку $x$:
\[
f(x)=
\begin{cases}
-x^2-x, & x<0,\\[2mm]
x^2-3x, & x\ge0.
\end{cases}
\]

\textbf{Левая часть, $x<0$:}
\[
y=-x^2-x.
\]
Ветви направлены вниз. Вершина:
\[
x_0=-\frac{b}{2a}=-\frac{-1}{2(-1)}=-\frac12,
\qquad
f\!\left(-\frac12\right)=\frac14.
\]
На допустимом промежутке нуль: $x=-1$.

\textbf{Правая часть, $x\ge0$:}
\[
y=x^2-3x.
\]
Ветви направлены вверх. Вершина:
\[
x_0=\frac32,
\qquad
f\!\left(\frac32\right)=-\frac94.
\]
Нули: $x=0$ и $x=3$.

\begin{center}
\begin{tikzpicture}
\begin{axis}[
  width=0.90\linewidth,
  height=0.50\linewidth,
  axis lines=middle,
  xmin=-3.5,xmax=4.2,
  ymin=-4.2,ymax=2.0,
  xtick={-3,-2,-1,0,1,2,3,4},
  ytick={-4,-3,-2,-1,0,1,2},
  xlabel={$x$},ylabel={$y$},
  grid=both,
  minor tick num=1,
  grid style={line width=.2pt,draw=gray!20},
  major grid style={line width=.3pt,draw=gray!30},
  samples=220,
  clip=true,
]
  \addplot[domain=-3.5:-0.001,very thick,color=accent] {-x^2-x};
  \addplot[domain=0:4.2,very thick,color=accent] {x^2-3*x};
  \addplot[domain=-3.5:4.2,dashed,color=warm] {0.25};
  \addplot[domain=-3.5:4.2,dashed,color=warm] {-2.25};
  \addplot[only marks,mark=*,mark size=2.3pt,color=darkaccent] coordinates {(-0.5,0.25) (1.5,-2.25)};
  \node[anchor=south west,color=darkaccent] at (axis cs:-0.5,0.25) {$\left(-\frac12;\frac14\right)$};
  \node[anchor=north west,color=darkaccent] at (axis cs:1.5,-2.25) {$\left(\frac32;-\frac94\right)$};
  \node[anchor=south east,color=warm] at (axis cs:4.1,0.25) {$m=\frac14$};
  \node[anchor=north east,color=warm] at (axis cs:4.1,-2.25) {$m=-\frac94$};
\end{axis}
\end{tikzpicture}
\end{center}

\subsection*{Параметр $y=m$: сколько общих точек?}
Горизонтальная прямая $y=m$ имеет с графиком ровно две общие точки в двух особых положениях --- когда касается одной из вершин и одновременно пересекает вторую ветвь:
\[
\boxed{m=\frac14\quad\text{или}\quad m=-\frac94}.
\]

\keyidea{Контрольная мысль:} недостаточно построить две полные параболы. Каждая формула действует только на своём промежутке: $x<0$ или $x\ge0$.

\clearpage

% ---------------- PIECEWISE + VECTORS ----------------
\sectiontitle{3. Кусочная функция: сначала смотрите на промежуток}

Для кусочной функции каждая формула работает только там, где указано условие рядом с ней.

Например, если гиперболическая формула содержит знаменатель $x$, но ветвь задана условием
\[
x<-4,
\]
то отдельная запись $x\ne0$ уже не добавляет нового ограничения: из $x<-4$ автоматически следует $x\ne0$.

\textbf{Проверяйте три вещи:}
\begin{itemize}
  \item границу промежутка: строгий знак $<$ или нестрогий $\le$;
  \item ОДЗ самой формулы;
  \item какие точки реально принадлежат конкретной ветви графика.
\end{itemize}

\vspace{4mm}
\sectiontitle{4. Векторы: стройте маршрут от начала к концу}

\keyidea{Правило сложения.} Чтобы сложить два вектора, совместите начало второго с концом первого. Вектор суммы соединяет начало первого с концом второго.

В параллелограмме $ABCD$ требуется выразить $\vect{D}{B}$ через $\vect{A}{D}$ и $\vect{A}{B}$. Идём из $D$ в $B$ через точку $A$:
\[
\vect{D}{B}=\vect{D}{A}+\vect{A}{B}.
\]
Но
\[
\vect{D}{A}=-\vect{A}{D}.
\]
Поэтому
\[
\boxed{\vect{D}{B}=\vect{A}{B}-\vect{A}{D}}.
\]

\begin{center}
\begin{tikzpicture}[scale=0.70]
  \tkzDefPoint(0,0){A}
  \tkzDefPoint(4.6,0){B}
  \tkzDefPoint(1.2,2.8){D}
  \tkzDefPoint(5.8,2.8){C}
  \tkzDrawPolygon[line width=0.9pt,color=softgray](A,B,C,D)
  \tkzLabelPoint[below left](A){$A$}
  \tkzLabelPoint[below right](B){$B$}
  \tkzLabelPoint[above left](D){$D$}
  \tkzLabelPoint[above right](C){$C$}
  \tkzDrawSegment[-{Stealth[length=3mm]},line width=1.5pt,color=accent](A,B)
  \tkzDrawSegment[-{Stealth[length=3mm]},line width=1.5pt,color=accent](A,D)
  \tkzDrawSegment[-{Stealth[length=3mm]},line width=1.6pt,color=warm](D,B)
  \node[color=accent,below] at (2.3,-0.15) {$\vect{A}{B}$};
  \node[color=accent,left] at (0.45,1.45) {$\vect{A}{D}$};
  \node[color=warm,right] at (3.55,1.35) {$\vect{D}{B}$};
\end{tikzpicture}
\end{center}

\textbf{Мини-памятка.} Направление важно: $\vect{D}{A}=-\vect{A}{D}$. Ищите маршрут от начала искомого вектора к его концу, записывайте векторы вдоль маршрута и учитывайте смену направления знаком «минус».

\clearpage

% ---------------- TRAINING 1 ----------------
\sectiontitle{5. Тренировочный блок — Путь А}

\textbf{10 заданий.} Сначала составьте модель или выполните построение самостоятельно. В задачах на движение следите за единицами и смыслом выбранной неизвестной.

\begin{enumerate}[label=\textbf{\arabic*.}]
  \item Два бегуна стартовали одновременно из одной точки круговой трассы в одном направлении. Скорость быстрого на $2$ км/ч больше скорости медленного. Быстрый завершил первый круг через $18$ минут. Через $20$ минут после старта медленному оставалось $500$ м до конца первого круга. Найдите скорость медленного бегуна.

  \item Первый лыжник проходит круг на $3$ минуты быстрее второго. За $60$ минут первый опережает второго ровно на один круг. За сколько минут второй лыжник проходит один круг?

  \item Первый спортсмен проходит круг на $4$ минуты быстрее второго. За $48$ минут первый опережает второго ровно на два круга. За сколько минут второй спортсмен проходит один круг?

  \item Второй бегун проходит круг за $x$ минут, первый --- на $2$ минуты быстрее. За $45$ минут первый опережает второго на один круг. Запишите уравнение относительно $x$ и укажите ограничение на $x$. Решать уравнение не требуется.

  \item Представьте функцию
  \[
  g(x)=|x|(x+2)-x
  \]
  в виде кусочной функции без знака модуля.
\end{enumerate}

\clearpage

% ---------------- TRAINING 2 ----------------
\sectiontitle{5. Тренировочный блок — продолжение}

\begin{enumerate}[label=\textbf{\arabic*.},start=6]
  \item Для функции
  \[
  f(x)=|x|(x-1)-2x
  \]
  найдите вершины обеих параболических частей и нули функции с учётом ограничений на $x$.

  \item Для той же функции $f(x)=|x|(x-1)-2x$ найдите все значения $m$, при которых прямая $y=m$ имеет с графиком ровно две общие точки.

  \item Дана кусочная функция, одна из ветвей которой задаётся формулой
  \[
  h(x)=\frac{4}{x}+2,\qquad x<-4.
  \]
  Нужно ли для этой ветви отдельно дописывать условие $x\ne0$? Кратко объясните.

  \item В параллелограмме $ABCD$ выразите вектор $\vect{D}{B}$ через векторы $\vect{A}{B}$ и $\vect{A}{D}$.

  \item В параллелограмме $ABCD$ пусть
  \[
  \vect{A}{B}=\vec a,\qquad \vect{A}{D}=\vec b.
  \]
  Выразите вектор $\vect{C}{A}$ через $\vec a$ и $\vec b$.
\end{enumerate}

\clearpage

% ---------------- ANSWERS ----------------
\sectiontitle{6. Ответы}

\begin{tabularx}{\linewidth}{@{}>{\bfseries}p{9mm}X@{}}
\toprule
№ & Ответ \\
\midrule
1 & $3$ км/ч. \\
2 & $15$ минут. \\
3 & $12$ минут. \\
4 & $\displaystyle \frac{45}{x-2}-\frac{45}{x}=1$, при $x>2$. \\
5 & $\displaystyle g(x)=\begin{cases}-x^2-3x,&x<0,\\x^2+x,&x\ge0.\end{cases}$ \\
6 & Левая часть: вершина $\left(-\frac12;\frac14\right)$, нуль $x=-1$. Правая часть: вершина $\left(\frac32;-\frac94\right)$, нули $x=0$ и $x=3$. \\
7 & $\displaystyle m=\frac14$ или $\displaystyle m=-\frac94$. \\
8 & Нет. Условие $x<-4$ уже гарантирует $x\ne0$. \\
9 & $\vect{D}{B}=\vect{A}{B}-\vect{A}{D}$. \\
10 & $\vect{C}{A}=-\vec a-\vec b$. \\
\bottomrule
\end{tabularx}

\vfill
{\small\color{softgray}Перед экзаменом особенно проверьте: единицы времени, ограничения на $x$, принадлежность точки нужной ветви и направление вектора.}

\end{document}
